% Options for packages loaded elsewhere
\PassOptionsToPackage{unicode}{hyperref}
\PassOptionsToPackage{hyphens}{url}
\PassOptionsToPackage{dvipsnames,svgnames,x11names}{xcolor}
\documentclass[
  10pt,
  fontset=fandol]{ctexart}
\usepackage{xcolor}
\usepackage[margin=16mm]{geometry}
\usepackage{amsmath,amssymb}
\setcounter{secnumdepth}{-\maxdimen} % remove section numbering
\usepackage{iftex}
\ifPDFTeX
  \usepackage[T1]{fontenc}
  \usepackage[utf8]{inputenc}
  \usepackage{textcomp} % provide euro and other symbols
\else % if luatex or xetex
  \usepackage{unicode-math} % this also loads fontspec
  \defaultfontfeatures{Scale=MatchLowercase}
  \defaultfontfeatures[\rmfamily]{Ligatures=TeX,Scale=1}
\fi
\usepackage{lmodern}
\ifPDFTeX\else
  % xetex/luatex font selection
\fi
% Use upquote if available, for straight quotes in verbatim environments
\IfFileExists{upquote.sty}{\usepackage{upquote}}{}
\IfFileExists{microtype.sty}{% use microtype if available
  \usepackage[]{microtype}
  \UseMicrotypeSet[protrusion]{basicmath} % disable protrusion for tt fonts
}{}
\makeatletter
\@ifundefined{KOMAClassName}{% if non-KOMA class
  \IfFileExists{parskip.sty}{%
    \usepackage{parskip}
  }{% else
    \setlength{\parindent}{0pt}
    \setlength{\parskip}{6pt plus 2pt minus 1pt}}
}{% if KOMA class
  \KOMAoptions{parskip=half}}
\makeatother
\setlength{\emergencystretch}{3em} % prevent overfull lines
\providecommand{\tightlist}{%
  \setlength{\itemsep}{0pt}\setlength{\parskip}{0pt}}
\usepackage{amsmath,amssymb,mathtools,needspace}
\xeCJKDeclareCharClass{CJK}{"2032,"0400->"04FF,"2160->"216B,"2460->"2473,"25A0,"25C7,"25CB,"25CF}
\IfFontExistsTF{Arial Unicode MS}{\xeCJKsetup{AutoFallBack=true}\setCJKfallbackfamilyfont{\CJKrmdefault}{Arial Unicode MS}}{}
\setcounter{secnumdepth}{0}
\setlength{\emergencystretch}{2em}
\allowdisplaybreaks
\usepackage{bookmark}
\IfFileExists{xurl.sty}{\usepackage{xurl}}{} % add URL line breaks if available
\urlstyle{same}
\hypersetup{
  pdftitle={重差加速的逻辑推理},
  colorlinks=true,
  linkcolor={Maroon},
  filecolor={Maroon},
  citecolor={Blue},
  urlcolor={Blue},
  pdfcreator={LaTeX via pandoc}}

\title{重差加速的逻辑推理}
\author{}
\date{}

\begin{document}
\maketitle

\subsubsection{12.4.3　重差加速的逻辑推理}\label{ux91cdux5deeux52a0ux901fux7684ux903bux8f91ux63a8ux7406}

假定所给逼近数列 \(\{a_n\}\) 由迭代公式

\[
a_{n+1}=\Phi(a_n)
\]

所生成，则有

\[
\begin{aligned}
a^*-a_{n+1}&=\Phi(a^*)-\Phi(a_n)\\
&=\Phi'(\xi_1)(a^*-a_n).
\end{aligned}
\]

又

\[
\begin{aligned}
a_{n+2}-a_{n+1}&=\Phi(a_{n+1})-\Phi(a_n)\\
&=\Phi'(\xi_2)(a_{n+1}-a_n),
\end{aligned}
\]

假定导数 \(\Phi'(\xi)\) 在所考察的范围内改变不大，则据上式得知

\[
\frac{a^*-a_{n+1}}{a^*-a_n}\approx\frac{a_{n+2}-a_{n+1}}{a_{n+1}-a_n}.
\]

这样，利用重差公设（12.4.1）可以断定

\[
\frac{a^*-a_{n+1}}{a^*-a_n}\approx\frac1\delta,
\]

即有

\[
a^*\approx\frac{\delta}{\delta-1}a_{n+1}-\frac1{\delta-1}a_n,
\]

因而有加速公式（12.4.2），重差加速法则得证。

上述推理过程自然会令人联想起经典数值分析中的松弛技术（参看第 4 章 4.3 节），两种加速公式的形式相同。

不过，更要区分传统的松弛技术与刘徽的重差技术两者的差异：

其一，松弛技术的前提是给出了具体的迭代函数 \(\Phi\)，并且要求函数 \(\Phi\) 具有一定的光滑性，而重差加速的理论分析所针对的迭代函数 \(\Phi\) 是虚构的，实际问题中并不存在。前述理论分析只是``借题发挥''而已。

其二，传统的松弛技术是近代学者的研究成果，而重差加速是中华先贤刘徽在两千年前提出的大智慧，两者不可同日而语。

其三，传统的松弛技术仅仅解决一些简单的数值分析例题，而刘徽的重差加速普遍适用于形形色色的大规模的计算工程，两者相比，真是``小巫见大巫''了。

值得指出的是，这里关于刘徽加速的讨论其实不能算作证明，充其量只是一种``说明''而已。逼近过程的收敛性与逼近加速的关系能否给出严格的逻辑证明，有待人们进一步深入研究。

\begin{quote}
校注（agent 补充）：原文``参看第 4 章 4.3 节''已保留。回看 PDF101（书页88），该节实际标题为``4.3 Romberg 算法''，此处对``松弛技术''的章节指引疑有误；未推定或代填作者意指的章节号。
\end{quote}

\subsection{本节覆盖原页的校注记录}\label{ux672cux8282ux8986ux76d6ux539fux9875ux7684ux6821ux6ce8ux8bb0ux5f55}

下列记录按原页关联，含同页相邻小节的校注；计算时同时读取上文校注和完整原页。

\begin{itemize}
\tightlist
\item
  \texttt{CH12-E005}：原书 PDF 页 306
\end{itemize}

\end{document}
